\chapter{System-on-Chip (SoC) Architecture and Interconnect}
\label{chap:soc_architecture}

\section{System Partitioning: Processing System vs. Programmable Logic}
The spGPU architecture adopts a hardware/software co-design strategy that leverages the Xilinx Zynq-7000 All-Programmable System-on-Chip (SoC). The system is physically partitioned into two tightly coupled domains:
\begin{enumerate}
    \item \textbf{Processing System (PS)}: Hosts the dual-core ARM Cortex-A9 processor running bare-metal C firmware. The CPU executes high-level tasks that require complex branching, floating-point calculations, and dynamic state evaluation, including rigid-body physics, ball-to-ball and boundary elastic collision resolution, and frame timing orchestration.
    \item \textbf{Programmable Logic (PL)}: Implements the spGPU coprocessor, containing the instruction dispatching unit, the parallel compute cores, the distributed tile memories, the performance analyzer, and the physical HDMI output engine.
\end{enumerate}

Figure~\ref{fig:soc_overview} illustrates the global block diagram of the SoC architecture, highlighting the interconnect interfaces, clock domains, and interrupt lines connecting the ARM processor with the FPGA logic.

\begin{figure}[htbp]
\centering
\resizebox{0.95\textwidth}{!}{
    \begin{tikzpicture}[
    >=Stealth,
    node distance=1.2cm and 1.5cm,
    every text node part/.style={align=center},
    box/.style={draw=black!80, thick, rounded corners=3pt, fill=white, inner sep=6pt, font=\small},
    psbox/.style={box, fill=blue!8, draw=blue!60},
    plbox/.style={box, fill=emerald!8, draw=emerald!60},
    corebox/.style={box, fill=orange!10, draw=orange!70, minimum width=2.2cm, minimum height=0.9cm},
    tilebox/.style={box, fill=purple!10, draw=purple!70, minimum width=2.2cm, minimum height=0.9cm},
    bus/.style={draw=black!70, line width=1.5pt, ->},
    sig/.style={draw=blue!70!black, thick, ->},
    intsig/.style={draw=red!80!black, thick, dashed, ->},
    subframe/.style={draw=gray!50, dashed, rounded corners=6pt, inner sep=10pt}
]

% ==================== PROCESSING SYSTEM (PS) ====================
\begin{scope}[local bounding box=ps_box]
    \node[psbox, font=\bfseries\normalsize, fill=blue!15, minimum width=13cm] (ps_title) {Processing System (PS) -- Dual-Core ARM Cortex-A9 (Zynq-7000)};
    
    \node[psbox, below=0.6cm of ps_title.south west, anchor=north west, minimum width=3.2cm, minimum height=1.6cm] (arm_app) {
        \textbf{Bare-Metal Application}\\
        \footnotesize 60 FPS Physics Engine\\
        \footnotesize Collision \& Game Loop
    };
    
    \node[psbox, right=0.8cm of arm_app, minimum width=3.0cm, minimum height=1.6cm] (arm_lib) {
        \textbf{spGPU Driver}\\
        \footnotesize\texttt{sp\_lib.c} / \texttt{sp\_lib.h}\\
        \footnotesize 64-bit Packet Builder
    };

    \node[psbox, right=0.8cm of arm_lib, minimum width=2.4cm, minimum height=1.6cm] (axi_dma) {
        \textbf{AXI DMA}\\
        \footnotesize Simple Mode MM2S\\
        \footnotesize Direct Memory Stream
    };

    \node[psbox, fill=red!10, draw=red!60, right=0.8cm of axi_dma, minimum width=2.2cm, minimum height=1.6cm] (gic) {
        \textbf{ARM GIC}\\
        \footnotesize Interrupt Controller\\
        \footnotesize IRQ 61 \& IRQ 62
    };

    \draw[bus] (arm_app) -- (arm_lib);
    \draw[bus] (arm_lib) -- (axi_dma);
\end{scope}
\node[subframe, fit=(ps_box), label={[anchor=north west, font=\bfseries\color{blue!80!black}]above left:PS (ARM Processing System Domain)}] {};

% ==================== INTERCONNECT ====================
\node[box, fill=gray!20, draw=gray!70, below=1.2cm of axi_dma, minimum width=3.8cm, font=\small\bfseries] (axi_stream) {
    AXI4-Stream (64-bit)\\
    \footnotesize\texttt{s\_axis\_tdata[63:0]}, \texttt{tvalid}, \texttt{tready}
};
\draw[bus] (axi_dma) -- (axi_stream);

% ==================== PROGRAMMABLE LOGIC (PL) ====================
\begin{scope}[local bounding box=pl_box]
    \node[plbox, fill=teal!15, font=\bfseries\normalsize, below=1.2cm of axi_stream, minimum width=15.5cm] (pl_title) {
        Programmable Logic (PL) -- Custom spGPU Multicore Accelerator (100 MHz Core Domain)
    };

    \node[plbox, below=0.7cm of pl_title.west, anchor=north west, minimum width=3.4cm, minimum height=1.2cm] (axi_wrap) {
        \textbf{AXI-Stream Wrapper}\\
        \footnotesize\texttt{myaxistream\_v1\_0}\\
        \footnotesize Protocol Conversion FSM
    };
    \draw[bus] (axi_stream.south) -| (axi_wrap.north);

    \node[plbox, fill=amber!15, draw=amber!70, right=0.8cm of axi_wrap, minimum width=4.0cm, minimum height=1.2cm] (sched) {
        \textbf{Dynamic Scheduler}\\
        \footnotesize\texttt{spScheduler} \& 10$\times$ FIFOs\\
        \footnotesize Bounding Box $[Y_{\min}, Y_{\max}]$
    };
    \draw[bus] (axi_wrap) -- node[above, font=\scriptsize] {\texttt{instr\_word[63:0]}} (sched);

    \node[plbox, fill=red!10, draw=red!60, right=1.2cm of sched, minimum width=3.2cm, minimum height=1.2cm] (analyzer) {
        \textbf{spANALYZER}\\
        \footnotesize 1-sec (100 MHz) Counter\\
        \footnotesize Exact FPS Measurement
    };

    % Core Array
    \node[corebox, below=1.2cm of sched.south west, anchor=north] (c0) {\textbf{spCORE 0}\\ \scriptsize Tile 0 ($Y: 0..23$)};
    \node[corebox, right=0.3cm of c0] (c1) {\textbf{spCORE 1}\\ \scriptsize Tile 1 ($Y: 24..47$)};
    \node[font=\bfseries, right=0.3cm of c1] (cdots) {\dots};
    \node[corebox, right=0.3cm of cdots] (c9) {\textbf{spCORE 9}\\ \scriptsize Tile 9 ($Y: 216..239$)};

    % Tile BRAM Array
    \node[tilebox, below=1.0cm of c0] (t0) {\textbf{Tile VRAM 0}\\ \scriptsize $320\times 24$ D-Buffer\\ \scriptsize + 4-bit Z-RAM};
    \node[tilebox, below=1.0cm of c1] (t1) {\textbf{Tile VRAM 1}\\ \scriptsize $320\times 24$ D-Buffer\\ \scriptsize + 4-bit Z-RAM};
    \node[font=\bfseries, below=1.3cm of cdots] (tdots) {\dots};
    \node[tilebox, below=1.0cm of c9] (t9) {\textbf{Tile VRAM 9}\\ \scriptsize $320\times 24$ D-Buffer\\ \scriptsize + 4-bit Z-RAM};

    \draw[bus] (sched.south) -- ++(0,-0.3) -| (c0.north);
    \draw[bus] (sched.south) -- ++(0,-0.3) -| (c1.north);
    \draw[bus] (sched.south) -- ++(0,-0.3) -| (c9.north);

    \draw[bus] (c0) -- node[right, font=\tiny] {$(X,Y,C,Z)$} (t0);
    \draw[bus] (c1) -- node[right, font=\tiny] {$(X,Y,C,Z)$} (t1);
    \draw[bus] (c9) -- node[right, font=\tiny] {$(X,Y,C,Z)$} (t9);

    % Scanline Multiplexer
    \node[box, fill=cyan!15, draw=cyan!70, below=0.8cm of tdots, minimum width=11cm, minimum height=0.9cm] (mux) {
        \textbf{Scanline Video Multiplexer}\\
        \footnotesize Selects Tile $0 \dots 9$ based on current vertical scanline $Y_{\text{scan}} \in [0, 239]$
    };
    \draw[bus] (t0.south) -- (t0.south |- mux.north);
    \draw[bus] (t1.south) -- (t1.south |- mux.north);
    \draw[bus] (t9.south) -- (t9.south |- mux.north);

    % Display Controller
    \node[box, fill=blue!10, draw=blue!60, below=0.8cm of mux.west, anchor=north west, minimum width=4.5cm] (vga) {
        \textbf{Display Controller (\texttt{vga\_struct})}\\
        \footnotesize VGA Timing: $640\times 480$ @ 60 Hz\\
        \footnotesize 2$\times$ Nearest-Neighbor Upscaler
    };
    \draw[bus] (mux.south -| vga.north) -- (vga.north);

    % TMDS & Serializer
    \node[box, fill=indigo!10, draw=indigo!60, right=0.8cm of vga, minimum width=3.8cm] (tmds) {
        \textbf{HDMI TMDS TX}\\
        \footnotesize TMDS 8b/10b Encoders\\
        \footnotesize\texttt{OSERDESE2} 5:1 DDR (125 MHz)
    };
    \draw[bus] (vga) -- node[above, font=\scriptsize] {RGB888} (tmds);

    \node[box, fill=slate!20, draw=black, right=0.8cm of tmds, minimum width=2.4cm] (hdmi_port) {
        \textbf{HDMI Port}\\
        \footnotesize Direct TMDS\\
        \footnotesize Differential Pairs
    };
    \draw[bus] (tmds) -- node[above, font=\tiny] {\texttt{CLK}, \texttt{D[2:0]}} (hdmi_port);

\end{scope}
\node[subframe, fit=(pl_box), label={[anchor=north west, font=\bfseries\color{teal!80!black}]above left:PL (FPGA Programmable Logic Domain)}] {};

% ==================== INTERRUPT CONNECTIONS ====================
\draw[intsig] (analyzer.north) to[out=90, in=0] node[pos=0.6, right, font=\scriptsize] {IRQ 61: 1 Hz Tick (FPS)} (gic.east);
\draw[intsig] (vga.west) to[out=180, in=-90] ++(-0.6, 3.5) to[out=90, in=180] node[pos=0.85, above, font=\scriptsize] {IRQ 62: 60 Hz VSYNC} (gic.south);

\end{tikzpicture}

}
\caption{spGPU System-on-Chip (SoC) Top-Level Architecture and Interconnect.}
\label{fig:soc_overview}
\end{figure}

\section{AXI DMA Communication Subsystem}
To offload graphics commands without burdening the CPU with per-pixel or per-word Programmed I/O (PIO) transactions, communication from the PS to the PL is mediated by the \textbf{Xilinx AXI Direct Memory Access (AXI DMA)} core configured in \textit{Simple Direct Register Mode} for Memory-Mapped to Stream (MM2S) transfers:
\begin{itemize}
    \item \textbf{High-Speed Memory Read}: The DMA controller accesses DDR memory over the 32-bit AXI High-Performance (HP) slave port on the PS.
    \item \textbf{Stream Conversion}: Instructions formatted into contiguous memory buffers are streamed over a 64-bit AXI4-Stream master interface (\texttt{M\_AXIS\_MM2S}).
    \item \textbf{Low Software Overhead}: The CPU merely writes the source address (\texttt{MM2S\_SA}) and transfer byte length (\texttt{MM2S\_LENGTH}) to the DMA memory-mapped control registers, allowing the hardware engine to burst instructions autonomously.
\end{itemize}

\section{AXI4-Stream Custom Peripheral (\texttt{myaxistream\_1\_0})}
The raw AXI4-Stream interface from the DMA engine is received on the FPGA by a custom conversion peripheral, \texttt{myaxistream\_1\_0} (detailed in \texttt{myaxistream\_v1\_0\_S\_AXIS.vhd}). This module serves as the primary gateway and protocol adapter between the Xilinx AXI bus and the internal scheduler of spGPU.

\subsection{State Machine and Handshaking}
The stream interface operates under a three-state Finite State Machine (FSM):
\begin{enumerate}
    \item \texttt{WAIT\_S}: The receiver remains in idle until the GPU scheduler asserts \texttt{ready\_enb} (indicating that the core FIFOs are ready to absorb a new command). Upon detecting a rising edge on \texttt{ready\_enb}, the FSM transitions to \texttt{TX1\_S}.
    \item \texttt{TX1\_S}: The core asserts \texttt{S\_AXIS\_TREADY <= '1'}, signalling to the DMA master that it is prepared to accept data. When \texttt{S\_AXIS\_TVALID} is sampled high, the FSM advances to \texttt{TX2\_S}.
    \item \texttt{TX2\_S}: The 64-bit data bus \texttt{S\_AXIS\_TDATA} is latched into \texttt{instr\_word}, \texttt{instr\_valid} is pulsed high for one cycle, \texttt{S\_AXIS\_TREADY} is deasserted, and the FSM returns to \texttt{WAIT\_S}.
\end{enumerate}

Listing~\ref{lst:axis_fsm} shows the core VHDL FSM implementation of this handshake protocol.

\begin{lstlisting}[language=VHDL, caption={AXI4-Stream Handshake and State Machine in \texttt{myaxistream\_v1\_0\_S\_AXIS.vhd}.}, label={lst:axis_fsm}]
process(all) 
begin
    if S_AXIS_ARESETN = '0' then
        next_state_t <= WAIT_S;
    else 
        case state_t is
            when WAIT_S => 
                if fsm_start_edge = '1' then
                    next_state_t <= TX1_S;
                else 
                    next_state_t <= WAIT_S;
                end if;
            when TX1_S =>
                if S_AXIS_TVALID = '1' then
                    next_state_t <= TX2_S;
                else 
                    next_state_t <= TX1_S;
                end if;
            when TX2_S =>
                next_state_t <= WAIT_S;
            when others => 
                next_state_t <= WAIT_S;        
        end case;
    end if;        
end process;
\end{lstlisting}

This handshake guarantees that incoming 64-bit instruction words from the DMA engine are ingested synchronously into the GPU clock domain without risk of FIFO overflow or lost words.
